\section{Conclusion}
\label{sec:conclusion}

New York taxi payment screens compute suggested tips as percentages of a
total that includes every government surcharge on the meter. We show
that this design choice has a measurable fiscal consequence: a
statutory dollar of surcharge raises the average credit-card tip by
about eleven cents---about fifteen cents per surcharge dollar the meter
actually records.

The estimate is well identified for a behavioral parameter of this kind.
It comes from surcharges that switch on a clock, assigned to the minute
of pickup, replicated at statutory doses from \$0.50 to \$5.00, in both
directions, at three separate times of day, against two independent
control series (weekends and exempt legal holidays), on a flat fare
where the confound cannot exist by construction (JFK),
and---crucially---reproduced by a design with no clock in it at all,
the spatial boundary of the 2025 congestion toll. Across six within-day changes the effect spans $0.090$ to
$0.117$ per statutory dollar ($\rho^{IV}$ of $0.122$--$0.150$ per
recorded dollar); the spatial estimate is $0.107$ per statutory dollar
against the clock designs' pooled $0.108$. Bandwidth and donut
variations move it by less than a cent. Two estimates, both at 6:00am on weekdays, fail, and
Section~\ref{sec:results-6am} shows why: the weekend estimates at that
same clock and dose are normal, so what breaks the weekday window is the
dawn compositional shift, not the mechanism.

Getting there required one correction we want to flag for anyone
replicating this design. The obvious specification---regress the tip on
an indicator for being past the cutoff---is badly biased, because the
metered fare jumps at 4:00pm along with the surcharge. Rush-hour trips
are slower and dearer, and part of the observed tip step is that
fare increase rather than the surcharge. Conditioning on the
non-surcharge base removes it, and the weekend placebo falls by roughly
three-quarters when it does. Designs that use time-of-day
discontinuities in prices should expect the same problem wherever the
underlying price moves with the time of day.

The interpretation is straightforward. Riders tip on a surcharge dollar
at essentially the rate they tip on a fare dollar, because the interface
gives them no way to distinguish the two. They do not switch to cash,
stop tipping, or abandon the default buttons in economically meaningful
numbers. What looks like generosity toward drivers is, at the margin,
arithmetic performed by a payment terminal.

The demand margins sharpen the picture into a single sentence: riders
do not retime hails around the 4:00pm fee, do not ride less when the
congestion toll arrives, do not switch to cash, and offset only a fifth
of a 150 percent surcharge increase---the fee is absorbed on every
margin except the one the software moves for them. It also yields a
concrete policy dial. One 2012 vendor already computed suggestions on
the fare alone, so the counterfactual is not hypothetical: mandating a
fare-only tip base would return roughly \$18.7 million a year of the 2019
surcharge's induced gratuities and \$2.2 million of the toll's to riders,
at zero collection cost, by changing one line of payment-terminal
arithmetic---about thirty cents per surcharge-paying card trip.

Two implications follow. For fiscal policy, the incidence of a fee
levied on a metered transaction does not stop at the statutory payment:
part of it continues past the agency to the worker, unlegislated and
uncounted---roughly \$18.7 million from the 2019 congestion surcharge and
\$2.2 million a year from the 2025 toll. For the growing literature on
interface-mediated payments, the result identifies a channel by which
policy and choice architecture interact: what a default is computed
\emph{on} matters as much as what it is set \emph{to}.

Three limitations bound the claims. Tips are observed only for card
payments, so the pass-through we estimate is conditional on card use,
which the design shows does not itself respond at the cutoff. Our
estimates aggregate across payment vendors whose adherence rates differ
by a factor of four, so $\rho$ is a market average rather than a
behavioral primitive. And the 2019 congestion-surcharge window predates
our minute-cell sample and abuts the pandemic; we treat it as
corroborative. Appendix~\ref{app:replication} ties the modern mechanism to the
canonical result: replicating \textcite{haggag2014default} on 2012--2013
data shows the two payment vendors of that era disagreed about the base
itself---one stopped at the fare plus surcharge, the other went on to
include tax and tolls---and that this base difference accounts for about
three-quarters of the vendor tip gap they documented. By 2025 every
vendor had converged on the inclusive base, which is what turns a
vendor-level quirk into a market-wide fiscal channel.
